\documentclass[10pt]{article}
\usepackage[margin=20mm]{geometry}
\usepackage{amsmath,amssymb,fontspec,unicode-math,xurl,hyperref}
\IfFontExistsTF{DejaVu Serif}{\setmainfont{DejaVu Serif}}{\setmainfont{Latin Modern Roman}}
\setmathfont{latinmodern-math.otf}
\setlength{\parindent}{0pt}\setlength{\parskip}{6pt}\sloppy
\begin{document}
Abstract: The published dynamic allocation rule can assign more than 100 percent pruning to a block, even with positive entropy scores and an exactly satisfied compensation budget. An explicit example follows every grouping and search step; this is a Major defect in the central allocation method as specified. A separate benchmark-average discrepancy is Minor. Published model scores are not reproduced or refuted by the allocation example, and the released evaluation snapshot cannot resolve the missing pruning implementation.\par
Date: Initial review, 24 October 2026\par
Title: Scientific review of multi-structure LLM pruning\par
\# Identity and inspected scope\par
Paper 0116, I25-43485: M. Yang, S. Lin, C. Li, and X. Chang, Let LLM Tell What to Prune and How Much to Prune, ICML 2025, PMLR 267, pp. 70833–70849 [Retained latex: \textbackslash{}cite\{paper\}]. The supplied original contains 17 pages; page numbers below are its printed pages 1–17.\par
The main method (Sections 3.1–3.4), experiments (Section 4), and Appendices A–C were inspected, including block grouping, compensation, dynamic ratios, row/column scoring, attention-group consistency, and ablations. All four files in the pinned public tree [Retained latex: \textbackslash{}cite\{code\}] were read statically. They contain evaluation/data-loading code and a README, but no implementation of the proposed pruning procedure. The README explicitly attributes its released models to reruns on different hardware and acknowledges seed-dependent differences from the paper. This is a recorded provenance qualification, not evidence that either set of reported results is false.\par
The review uses independent exact arithmetic, not author-code execution or new model experiments. Benchmark training/evaluation, latency measurements, external model downloads and statistical reruns were not performed. The original and supplied-file hashes were checked.\par
\# Definitions and checked reasoning\par
For \(N\) blocks let \(s_i\) denote the entropy difference in Eq. (1), p. 3. Equation (4), p. 4, divides scores at \(\mu_s\pm\alpha\sigma_s\) into \(S_L,S_M,S_H\). High-score blocks are preserved, medium-score blocks receive \(p_{\rm low}\), and Eq. (5) calculates their leftover pruning load. Candidate full-block removals \(F\) are the \(|S_H|\) lowest-score members of \(S_L\); candidate layer removals \(E\) depend on \(S_M\). The depth-first search minimizes excess compensation subject to covering the leftover, preferring more full-block removals among ties.\par
Equation reference: eq:allocation. For \(G\) remaining blocks, Eq. (8), p. 5, assigns \[p_g=\frac{G\omega_g}{\sum_h\omega_h}\widehat p,\qquad
 \omega_g=(\bar s/s_g)^\beta,\qquad
 \widehat p=\frac{G p_{\rm target}-p_{\rm compensation}^{*}+p_{\rm leftover}}{G}.
 \] For positive weights this preserves the sum \(\sum_g p_g=G\widehat p\). It does not by itself enforce the necessary individual constraint \(0\leq p_g\leq1\). The following example tests precisely that missing implication.\par
The Gaussian entropy formula’s final expression, \(H=\log\sigma+\tfrac12\log(2\pi)+\tfrac12\), and Figure 4’s approximate scores were independently checked. No substantive defect was found in that finite example. The output-entropy difference is treated here as the paper’s operational heuristic; no general causal or rank-preservation theorem is assumed.\par
\# Supported problems\par
Problem 1 (Major: a reachable allocation exceeds the feasible pruning range).  Claim and location. Section 3.2, Eqs. (4)–(8), pp. 4–5, and Algorithm 1 purport to assign a pruning ratio to each block while meeting the target budget. Dynamic allocation is a named central contribution.\par
Parameters and assumptions. Take eight equal-size blocks with scores \[(10^{-6},2\cdot10^{-6},3\cdot10^{-6},0.1,0.1,2,2,2).\] Set \(\alpha=1/2\), \(\beta=1\), \(p_{\rm target}=0.4\) and \(p_{\rm low}=0.2\). All scores and both scaling parameters are positive; \(0<p_{\rm low}<p_{\rm target}<1\). The paper supplies no further restriction excluding these values. The scores are compatible with the stated Gaussian entropy calculation: output standard deviation \(1\) and masked standard deviations \(e^{-s_i}\) give exactly the designated differences. This is an allocation-level counterexample on admissible score inputs, not a claim to have measured these scores in the published LLaMA checkpoints.\par
Complete path through the algorithm. With population standard deviation as specified in Eq. (4), \[\mu_s=0.77500075,\quad \sigma_s=0.949670384,\quad
 T_l=0.300165558,\quad T_h=1.249835942.\] Consequently \(S_L\) contains the first five blocks, \(S_H\) the last three, and \(S_M\) is empty. Thus \(E\) is empty, \(F\) consists of the first three blocks, and \[p_{\rm leftover}=3(0.4)=1.2.\] Each full-block removal contributes \(1-0.4=0.6\) in Eq. (7). Every optimal search solution removes exactly two members of \(F\), giving compensation exactly \(1.2\) and excess zero. Removing one is infeasible; removing three has positive excess. Thus neither tie-breaking nor an approximation in the search can evade the case.\par
Equation reference: eq:invalid. The remaining set has scores \((\varepsilon,0.1,0.1)\), where \(\varepsilon\in\{10^{-6},2\cdot10^{-6},3\cdot10^{-6}\}\), and \(G=3\), \(\widehat p=0.4\). Cancelling the common \(\bar s\) factor in the weights gives \[p_\varepsilon=\frac{1.2}{1+20\varepsilon}>1,\qquad
 p_{0.1}=\frac{12\varepsilon}{1+20\varepsilon}.
 \] For \(\varepsilon=3\cdot10^{-6}\) these are \((1.1999280043,0.0000359978,0.0000359978)\). The total over all eight blocks is the intended \(3.2\), but a block cannot lose \(119.99\%\) of its parameters. The independent script enumerates all three optimal search solutions; each has the same violation. This example avoids zero scores, negative scores, empty remaining sets, infeasible compensation and unequal block sizes.\par
Alternative readings and impact. A hidden clipping or redistribution step could make an implementation usable. No such step occurs in the published formula or pseudocode, and the pinned repository contains no pruning implementation to establish one. Clipping alone breaks the intended total: the displayed case then removes about \(3.000072\) block-equivalents, or \(37.5009\%\), instead of \(40\%\). Actual experimental scores need not have entered this regime. The finding establishes a defect in the core algorithmic specification, not a benchmark collapse or a claim that the released checkpoints contain impossible sparsity.\par
Precisely bounded correction. For positive weights and a feasible residual budget \(B=G\widehat p\in[0,G]\), choose \[p_g=\min\{1,\lambda\omega_g\},\qquad \sum_g p_g=B.\] The sum is continuous, nondecreasing from \(0\) to \(G\), so a solution exists; it is unique for \(0<B<G\). Here it yields \((1,0.1,0.1)\) and restores the total \(3.2\). This verifies box feasibility and budget conservation only. It does not prove optimal accuracy, reproduce the authors’ implementation, handle nonpositive entropy scores, or account for indivisible heads and actual layer parameter costs. Those require additional specification.\par
Severity and evidence. Major: Eqs. (4)–(7) feed Eq. (8), which feeds the within-block pruning stage. A valid ratio is essential to the proposed dynamic method. Evidence is the original p. 5 image, the exact arithmetic script, and its complete output in `independent\_checks.json`.\par
Problem 2 (Minor: a baseline’s reported average understates its displayed performance).  Claim and location. Appendix B, Table 4, p. 14, gives SliceGPT at \(40\%\) sparsity on LLaMA2-13B the five scores \((60.69,59.30,42.50,56.44,34.13)\) and average \(47.30\).\par
Check. The equally weighted mean is \[(60.69+59.30+42.50+56.44+34.13)/5=50.612.\] Rounding the displayed component scores cannot explain a \(3.312\)-point difference. The released evaluation script’s `calculate\_avg\_accuracy` also uses an unweighted task mean, although this later script is not proof of the historical computation.\par
Alternatives and bounded impact. A copied aggregate or a mistaken component cell could explain the discrepancy; raw result records are required to establish which. If the five component scores are correct, the aggregate is approximately \(50.61\), and the reported method’s advantage over this baseline falls from approximately \(8.74\) to \(5.44\) points. Its mean remains higher, so this does not overturn the within-table winner. The issue is Minor because it overstates a numerical comparison without demonstrated damage to the central comparative conclusion. The correction is conditional on the component scores, not a reconstruction of unobserved runs. Original p. 14 was visually checked to exclude an extraction or row-alignment error.\par
\# Evidence limits, counterevidence, and remaining work\par
Dynamic-ratio evidence. Table 13 reports static and dynamic means \(45.45\) and \(45.52\). The displayed components reproduce a \(0.068\) percentage-point difference; some tasks improve and others decline. This is an observed single-table difference, not established statistical robustness. The paper’s released-model README expressly notes seed sensitivity. Replicated pruning seeds and matched uncertainty estimates are needed to quantify the incremental dynamic-ratio benefit. This uncertainty does not supply the evidence for P1; the algebraic counterexample does.\par
Implementation boundaries. Actual \(\alpha,\beta,p_{\rm low}\) choices, clipping/renormalization, clustering counts, discrete head rounding, and handling of zero or signed entropy differences were not recovered from pruning source. The compensation formula for partial layer removal also needs actual parameter-count accounting when attention and MLP sizes differ. These are Needs Evidence, not additional demonstrated benchmark errors.\par
The bias-compensation notation in Eqs. (12)–(13) requires a dimensionally explicit mask. A column-mask vector with broadcasting can make the intended expression coherent; a full weight-shaped mask would instead suggest \(((1-M)\odot W)\widehat X\). That alternative prevents treating the notation alone as a demonstrated implemented defect. Weighting by activation variability is a heuristic, not a proven optimum asserted by this review.\par
No issue found within a specific code check. The evaluation’s multiplication of per-batch mean loss by sequence length \(S\), followed by division by the total number of sequences times \(S\), cancels. Since each scored sequence has the same \(S-1\) prediction positions, this correctly averages those mean losses; an apparent off-by-one issue was rejected. WikiText2 training and test splits are separate in the read loader. This does not certify preprocessing equivalence across every baseline.\par
Versions and missing evidence. A bounded official-publication/repository search did not identify a verified corrective version for these findings. This is not an exhaustive erratum claim. The public source is a later evaluation release, with the rerun caveat preserved. Separate supplement availability remains undetermined after the supplied OpenReview 403 access failure. No author response, adoption or human assessment is present.\par
Item | Judgment | Verified correction | Remaining work\par
P1 | Major | Capped allocation preserves feasible continuous budget | Historical algorithm, discrete costs, performance\par
P2 | Minor | \(50.612\) mean of listed cells | Original per-task records\par
Other | Needs Evidence | None claimed | Settings, seeds, latency and provenance\par
Overall supported severity is Major. Scientific decisions remain pending and human reviews are empty. This is an initial review, not acceptance; no author adoption has been requested.\par
\# Reproducibility and provenance\par
The supplied original SHA-256 is\par
> `79a75cab0ea573b3f09ee7d7c39001fce959eb07eb809411299233c755d59072`.\par
Manifest checks, extracted original text, the original p. 5 and p. 14 images, line-numbered copies of all four source files, source hashes, the reviewer-written Python script and its exact/finite outputs reside in `verification/0116/evidence/`. No supplied original or source was modified. The initial check read occurred before text extraction completed; the subsequent extraction succeeded and the actual text was read. All substantive calculations are independent of author-code execution.\par
The report was compiled twice with the preinstalled compiler, using\par
> pdflatex -no-shell-escape -interaction=nonstopmode
> -halt-on-error 0116.tex\par
All final PDF pages were rendered and actually inspected. The source and final binary hashes, compiler outputs, transfer verification and visual inspection receipt are saved alongside the coverage, findings and provenance records. No software was installed. Temporary compilation occurred in the available isolated runtime; final durable artifacts are in the assigned staging location.\par
The current assigned AI conversation performed the scientific work and visual inspection. No human scientific review is recorded. Task: [internal task]; batch token `[internal identifier omitted]`.\par
9 M. Yang, S. Lin, C. Li, and X. Chang, Let LLM Tell What to Prune and How Much to Prune, PMLR 267 (2025), 70833–70849. https://proceedings.mlr.press/v267/yang25p.html. Supplied 17-page proceedings original. PruneLLM, supplied commit `a5930e2737216179882ea7af4c0f142cdc0dfa02`, acquired 24 October 2026. https://github.com/yangmzevery/PruneLLM. Files inspected: `README.md`, `eval\_prune\_llm.py`, `eval.py`, `data.py`.\par
\end{document}
